\chapter{比例}\label{ch:g7:prop}

六年级学过比例表（\cref{ch:g6:propdata}）；本章将比例发展为
一个完整的工具：用交叉法则求第四个数值、
百分数，以及地图和模型的比例尺。
速度是最著名的比例关系，
将在\cref{ch:g8:speed}中学习。

\section{识别与补全比例表}

\begin{definition}[比例表]\label{def:g7:prop:table}
若一张两行的表格，其第二行的数都是第一行对应的数乘以同一个固定的数，则称该表为\emph{比例
表}\index{比例表}，此固定的数称为
\emph{系数}。
等价地，所有列的商
$\frac{\text{下}}{\text{上}}$ 都相等。
\end{definition}

\begin{example}\label{ex:g7:prop:recognize}
\begin{center}
\renewcommand{\arraystretch}{1.3}
\begin{tabular}{l|ccc}
上 & $4$ & $10$ & $14$ \\
\hline
下 & $6$ & $15$ & $21$ \\
\end{tabular}
\end{center}
商：$\frac64 = 1.5$，$\frac{15}{10} = 1.5$，
$\frac{21}{14} = 1.5$：成比例，系数为 $1.5$。
\end{example}

\begin{theorem}[交叉法则]\label{thm:g7:prop:cross}
在比例表中，列为 $\begin{pmatrix} a \\ b
\end{pmatrix}$ 与 $\begin{pmatrix} c \\ d \end{pmatrix}$ 的两列，其交叉积
相等：
\[
a \times d = b \times c .
\]
因此，由其余三个数可以算出第四个数：
$d = \dfrac{b \times c}{a}$。
\end{theorem}

\begin{proof}
设 $k$ 为系数：$b = ka$，$d = kc$。于是
\[
a \times d = a \times kc = k \times ac,
\qquad
b \times c = ka \times c = k \times ac :
\]
两个交叉积都等于 $k \times ac$。将 $ad = bc$ 两边除以 $a$，
即得 $d$ 的公式。
\end{proof}

\begin{example}[求第四个比例项]\label{ex:g7:prop:fourth}
若 $7$ 本相同的书重 $2.8$ 千克，则 $12$ 本这样的书重多少？
列表并用交叉法则：
\begin{center}
\renewcommand{\arraystretch}{1.3}
\begin{tabular}{l|cc}
本数 & $7$ & $12$ \\
\hline
千克 & $2.8$ & $m$ \\
\end{tabular}
\end{center}
\[
7 \times m = 2.8 \times 12
\quad\Longrightarrow\quad
m = \frac{2.8 \times 12}{7} = \frac{33.6}{7} = 4.8 \text{ 千克}.
\]
按单位检验：一本书重 $0.4$ 千克，十二本重 $4.8$ 千克。
\end{example}

\section{百分数}

\begin{method}[百分数的应用与求法]\label{met:g7:prop:percent}
\begin{enumerate}
  \item \emph{求}$t\,\%$：乘以 $\frac{t}{100}$
        （$12\,\%$ 作用于 $350$，得 $0.12 \times 350 = 42$）；
  \item \emph{求}某部分占整体的百分数：计算
        $\frac{\text{部分}}{\text{整体}}$，再化为以 $100$ 为分母的分数
        （$21$ 占 $60$ 的比例：$\frac{21}{60} = \frac{35}{100} = 35\,\%$）；
  \item 总要问：\emph{这是谁的百分之几？}
        参照的整体很重要。
\end{enumerate}
\end{method}

\begin{example}\label{ex:g7:prop:percent}
某校的女生有 $180$ 名，占全校 $400$ 名学生。百分数为：
\[
\frac{180}{400} = \frac{180 \div 4}{400 \div 4} = \frac{45}{100}
= 45\,\% .
\]
反过来：$55\,\%$ 的学生是男生，男生有多少人？
$0.55 \times 400 = 220$。检验：$180 + 220 = 400$。
\end{example}

\section{比例尺}

\begin{definition}[比例尺]\label{def:g7:prop:scale}
若地图或模型的\emph{比例尺}\index{比例尺}为
$\frac{1}{n}$（记作 $1 : n$），则图上距离与实际距离成正比例，
系数为 $\frac1n$：
\[
\text{图上距离} = \frac{1}{n} \times \text{实际距离},
\]
两者须用\emph{同一单位}度量。
\end{definition}

\begin{example}\label{ex:g7:prop:scale}
在比例尺为 $1 : 50\,000$ 的地图上，两镇相距 $7$ 厘米。实际距离为：
\[
7 \times 50\,000 = 350\,000 \text{ 厘米} = 3\,500 \text{ 米} = 3.5
\text{ 千米}.
\]
反过来，一条 $12$ 千米的徒步路线在地图上的长度为
$12$ 千米 $= 1\,200\,000$ 厘米，所以 $1\,200\,000 \div 50\,000 = 24$ 厘米。
\end{example}

\begin{omfigure}
\begin{tikzpicture}
\begin{axis}[omaxis,
  width=8.4cm, height=5.6cm,
  xmin=0, xmax=8.6, ymin=0, ymax=34,
  xlabel={地图 (cm)}, ylabel={实际 (km)},
  xtick={2,4,6,8}, ytick={10,20,30}]
  \addplot[omDef, thick, domain=0:8.2] {4*x};
  \addplot[omDef, only marks, mark=*, mark size=1.5pt] coordinates
    {(2,8) (4,16) (6,24) (8,32)};
  \draw[dashed, omThm] (axis cs:7,0) -- (axis cs:7,28)
    -- (axis cs:0,28);
\end{axis}
\end{tikzpicture}

{\small 比例尺表示一种成正比例的关系：地图上的距离与实际距离对应，
在一条过原点的直线上（此处
$1$ 厘米 $\leftrightarrow$ $4$ 千米；虚线表示 $7$ 厘米
$\leftrightarrow$ $28$ 千米）。}
\end{omfigure}

\begin{remark}
图像是检验是否成正比例最快捷的方法
（见\cref{ch:g6:propdata}）：直线上的点若\emph{过原点}，
则所表示的两个量成正比例；不过原点的直线（如
出租车计费含固定起步价）表示\emph{不}成正比例，
尽管它也是直线。一次函数将在
\cref{ch:g9:linfunc}中使这一点更加精确。
\end{remark}

\section{练习}

\begin{exercise}[$\star$]\label{exo:g7:prop:1}
哪些表格是比例表？若比例系数存在，
请给出该系数。
\begin{center}
\renewcommand{\arraystretch}{1.2}
\begin{tabular}{l|ccc}
 & $3$ & $7$ & $11$ \\
\hline
 & $7.5$ & $17.5$ & $27.5$ \\
\end{tabular}
\qquad
\begin{tabular}{l|ccc}
 & $2$ & $5$ & $9$ \\
\hline
 & $6$ & $15$ & $28$ \\
\end{tabular}
\end{center}
\end{exercise}

\begin{exercise}[$\star$]\label{exo:g7:prop:2}
用交叉相乘法补全，并写出计算过程：
$5$ kg 土豆售价 $6$ 欧元；买 $8$ kg 需要多少欧元？
\end{exercise}

\begin{exercise}[$\star$]\label{exo:g7:prop:3}
一台打印机打印 $36$ 页用时 $3$ 分钟。按相同速度，
$5$ 分钟可打印多少页？打印 $96$ 页需要多少分钟？
\end{exercise}

\begin{exercise}[$\star$]\label{exo:g7:prop:4}
计算：$30\,\%$ 的 $250$；\quad $15\,\%$ 的 $60$；\quad 求 $t\,\%$，
使 $t\,\%$ 的 $80$ 为 $20$。
\end{exercise}

\begin{exercise}[$\star$]\label{exo:g7:prop:5}
一名篮球运动员投篮 $25$ 次，命中 $18$ 次。她的
命中率是多少？（把这个分数化为分母为 $100$ 的分数。）
\end{exercise}

\begin{exercise}[$\star$]\label{exo:g7:prop:6}
在比例尺为 $1 : 25\,000$ 的地图上，一条小路长 $9$ 厘米。它的实际
长度是多少千米？一个湖长 $2$ 千米，它在地图上长多少厘米？
\end{exercise}

\begin{exercise}[$\star$]\label{exo:g7:prop:7}
一辆模型汽车按比例尺 $1 : 43$ 制作。真实汽车长 $4.3$ 米。
模型长多少厘米？
\end{exercise}

\begin{exercise}[$\star\star$]\label{exo:g7:prop:8}
一份食谱用 $240$ 克巧克力制作 $8$ 人份。爱丽丝
有 $300$ 克巧克力，按整数份计，这些巧克力够做多少
人份？
\end{exercise}

\begin{exercise}[$\star\star$]\label{exo:g7:prop:9}
一件夹克的价格从 $80$ 欧元降到 $60$ 欧元。
\begin{enumerate}
  \item 降价多少欧元？降价占原价的
        百分之几？
  \item 之后价格从 $60$ 欧元涨回 $80$ 欧元。请解释
        为什么这次涨幅\emph{不}是 $25\,\%$ --- 算出正确的
        增长百分比。
\end{enumerate}
\end{exercise}

\begin{exercise}[$\star\star$]\label{exo:g7:prop:10}
表格中有两行成正比例：
\begin{center}
\renewcommand{\arraystretch}{1.2}
\begin{tabular}{l|ccc}
 & $4$ & $x$ & $10$ \\
\hline
 & $y$ & $10.5$ & $17.5$ \\
\end{tabular}
\end{center}
求比例系数，再求 $x$ 和 $y$。
\end{exercise}

\begin{exercise}[$\star\star\star$]\label{exo:g7:prop:11}
一台复印机把文件缩小为原尺寸的 $80\,\%$。取一条
$10$ cm 的线段先复印一次，再以相同
设置把\emph{副本}再复印一次。
\begin{enumerate}
  \item 复印一次后，线段长多少？复印两次后呢？
  \item 为什么复印两次后的结果不是原长的 $60\,\%$？
        两次复印等价于一次缩小为原长的百分之几？
\end{enumerate}
\end{exercise}

\section{问题：用一根木棍测量地球}

\begin{problem}[{周末问题 --- 影子成正比例，以及埃拉托色尼如何
于公元前 240 年算出地球的周长
}]
\label{pb:g7:prop:1}
二十二个世纪前，没有望远镜，没有卫星，也没有计算器，
亚历山大图书馆的馆长测量了地球 ---
借助一根木棍、一口井、一支骆驼商队，以及本章的数学。
本题重走他的足迹：先看影子的比例关系，
再完成那个著名的计算，
最后看看他的答案在人类尺度上是什么样子。

\medskip
\noindent\textbf{第一部分 --- 木棍的影子。}
太阳离我们如此遥远，以至于它射来的光线彼此
平行。在某一时刻，所有竖直物体与它们的影子
因此成正比例。
\begin{enumerate}
  \item 正午时，一根长 $1$ m 的竖直木棍影长 $0.4$ m，
        一棵树的影长为 $3.2$ m。
        这棵树有多高（见\cref{thm:g7:prop:cross}）？
  \item 同一时刻，一座塔的影长为 $26$ m。这座塔有多高？
        一个身高 $1.5$ m 的孩子，他的影子
        有多长？
  \item 用一句话说明，为什么这些计算
        只在一天中的\emph{同一时刻}才成立。
  \item 传说泰勒斯测量大金字塔时，等待的是这样
        一个时刻：\emph{自己的影子恰好与自身身高一样长}。
        此时的比例系数是多少？
        若金字塔影子的顶端与其塔尖正下方的地面点
        相距 $147$ m，
        金字塔有多高？
  \item 小结第一部分：在某一给定时刻，木棍、树、塔
        和金字塔的哪一个量是相同的，
        （见\cref{def:g7:prop:table}）？
\end{enumerate}

\medskip
\noindent\textbf{第二部分 --- 埃拉托色尼的计算。}
埃拉托色尼掌握两条事实。在塞恩城，
夏至正午，太阳\emph{恰好直射}：阳光
能照到最深的井底，竖直的木棍没有影子。
在正北方 $800$ km 处的亚历山大城，同一时刻，
一根竖直的木棍\emph{确实}投下了影子，太阳光线
与竖直方向所成的角为 $7.2^\circ$。
\begin{enumerate}[resume]
  \item 说明这两条观察结果为何共同证明：
        塞恩与亚历山大两地的地面不平行
        --- 也就是说，地球表面是弯曲的。
        （若地球是平的，平行光线照射两根木棍，
        会出现什么情况？）
  \item 在一幅画有平行太阳光线的圆形地球图上，
        亚历山大处的 $7.2^\circ$ 在地心处重新出现，
        即两座城市所在方向之间的夹角。
        请动手作图，自己验证一下（严谨的论证
        所用的等角对，见
        \cref{ch:g7:triangles}）。
  \item $7.2^\circ$ 是一个周角的几分之几？
  \item 从塞恩到亚历山大的弧长（$800$ km）对应
        $7.2^\circ$，整个周长对应
        $360^\circ$。列出比例式，并求
        地球的周长。
  \item 取 $\pi \approx
        3.14$（见\cref{prop:g6:measure:circle}），推出地球的直径，并四舍五入到
        最接近的百千米。（现代值为
        $12\,742$ km --- 公元前 240 年手持木棍的人，
        结果离它有多近？）
\end{enumerate}

\medskip
\noindent\textbf{第三部分 --- 人类尺度下的地球。}
\begin{enumerate}[resume]
  \item 按比例尺 $1 : 40\,000\,000$ 制作一个地球仪。利用
        第 9 题求得的周长，说明地球仪的周长恰好为
        $1$ m。它的直径是多少
        （精确到厘米）？
  \item 勃朗峰高约 $4.8$ km。把 $4.8$ km 换算为厘米，
        再除以 $40\,000\,000$，并用毫米表示
        这座山在地球仪上的高度。这个答案说明
        地球实际上有多“凹凸不平”？
  \item 一位步行者每天走 $40$ km。按这一速度，
        走完埃拉托色尼测得的整个周长需要多少天 ---
        大约相当于多少年？
  \item 可供呼吸的大气主要集中在地面以上
        约 $10$ km 的范围内。这一高度占地球半径
        （约 $6\,370$ km）的百分之几？
        （见\cref{met:g7:prop:percent}）。结果四舍五入到
        百分数的十分位。
  \item 史学家估计，埃拉托色尼宣布的数值
        换算为现代单位后，可能约为
        $42\,000$ km。取 $40\,000$ km 为真值，
        求他的误差百分数。用一句话总结
        关于木棍、水井与比例的道理。
\end{enumerate}
\end{problem}
